\subsection{Conclusión comparativa del eje transversal}
Los seis casos M6/M8 conservaron presión, capas y fijaciones idénticas. La comparación
no usa como referencia equivalente el M6 articulado de la revisión anterior: cambiar
empotramientos y añadir acoplamiento también modifica su respuesta.
Para la capa supuesta de $1\,\mathrm{mm}$, M8 reduce la flecha un $67.66\,\%$ y la tensión nominal un $56.38\,\%$ frente a M6. Para la capa supuesta de $2\,\mathrm{mm}$, M8 reduce la flecha un $68.53\,\%$ y la tensión nominal un $57.56\,\%$ frente a M6. Para la capa supuesta de $4\,\mathrm{mm}$, M8 reduce la flecha un $69.44\,\%$ y la tensión nominal un $58.65\,\%$ frente a M6. 
La mayor rigidez reduce la apertura de holgura y aumenta ligeramente el área cubierta y las cargas convergidas;
por ello no se supone que el cambio de varilla mantenga las fuerzas fijas.
El M8 conserva holgura positiva y menor que la capa supuesta en todos los casos, bajo los
empotramientos y propiedades de referencia. No mantiene exactamente los $0,5\,\mathrm{mm}$ iniciales.
Sin tolerancia funcional admisible, resistencia del producto y rigidez de uniones no se declara
que la separación sea adecuada para el funcionamiento real ni que el eje esté aprobado.
La siguiente definición de diseño debe confirmar clase/paso/raíz de la varilla, fijaciones enABS,
asientos608, tolerancias de holgura y validación del material compactado. No se modifica CAD.

El límite elástico tendría que superar al menos $315.36\,\mathrm{MPa}$ para compatibilidad nominal con la flexión elástica del caso M8 más exigente, aun sin concentradores ni reserva de diseño. Este umbral necesario no es un factor de seguridad ni certifica el material.
